% ============================================================
% CAPÍTULO 1 — Introducción y Contexto
% ============================================================
\chapter{Introducci\'on y Contexto}
\label{cap:introduccion}

\section{Proyecto Demeter --- s\'intesis}

Demeter nace en \textbf{diciembre de 2024} como iniciativa de la
Asociaci\'on Universitaria ESIBot para abordar la necesidad de
herramientas asequibles de monitorizaci\'on ambiental e investigaci\'on
agr\'icola de precisi\'on. El sistema, definido por su arquitectura de
cuatro capas, integra hardware embebido (ESP32), computaci\'on en el
borde (Raspberry Pi), infraestructura cloud (Google Cloud Platform) y
una interfaz web reactiva, articulados por un protocolo binario
propietario optimizado para nodos de bajo consumo.

El prop\'osito original no era construir una plataforma IoT gen\'erica,
sino resolver un problema concreto y bien acotado: en el
\textbf{Grado en Ingenier\'ia Agr\'icola} hay pr\'acticas docentes cuyas
plantas deben regarse y vigilarse a diario, y ese riego lo ven\'ian
haciendo a mano los t\'ecnicos del Departamento de Agronom\'ia. Demeter
deb\'ia \textbf{automatizar ese riego} y liberar as\'i horas de trabajo
manual repetitivo, aportando de paso la trazabilidad de cada planta y
cada experimento que un protocolo agron\'omico riguroso exige y que una
libreta de campo no da.

Conviene retener el dato temporal, porque enmarca todo el an\'alisis de
este informe: entre esa fecha de nacimiento y la redacci\'on del presente
documento median \textbf{veinte meses}, de los cuales el desarrollo
efectivo del sistema ocupa \textbf{cuarenta y cuatro d\'ias naturales}
(227 \textit{commits} entre enero y marzo de 2026). El resto del tiempo
el proyecto estuvo detenido, no por falta de voluntad, sino por la
ausencia de personas disponibles para empujarlo. Esa proporci\'on entre
tiempo transcurrido y tiempo trabajado es, en s\'i misma, la primera
evidencia del problema de sostenibilidad que el
cap\'itulo~\ref{cap:sostenibilidad} analiza en detalle.

La descripci\'on t\'ecnica completa del sistema se recoge en la
\textit{Documentaci\'on T\'ecnica del Sistema} (documento hermano de
84~p\'aginas). El presente informe se centra en el an\'alisis de
viabilidad y la propuesta de cierre del proyecto.

\section{Marco institucional}

ESIBot es una asociaci\'on universitaria con sede en la
\textbf{Escuela T\'ecnica Superior de Ingenier\'ia (ETSI)} de la
Universidad de Sevilla. Re\'une aproximadamente a \textbf{cien miembros}
y trabaja en cuatro l\'ineas: \textbf{rob\'otica}, \textbf{electr\'onica},
\textbf{telem\'atica} y \textbf{dise\~no 3D}. Su actividad se organiza en
proyectos internos de duraci\'on variable, en los que los miembros
participan de forma voluntaria y en paralelo a sus estudios.

El Proyecto Demeter se enmarc\'o como Proyecto Docente bajo el acuerdo
de colaboraci\'on con el \textbf{Departamento de Agronom\'ia de la
Escuela T\'ecnica Superior de Ingenier\'ia Agron\'omica (ETSIA)}. El
objetivo pactado era proveer un \textbf{sistema de riego automatizado}
que sustituyera el trabajo manual de los t\'ecnicos encargados de las
plantas asignadas a una pr\'actica docente del Grado en Ingenier\'ia
Agr\'icola, a\~nadiendo el registro y la trazabilidad de las variables
ambientales de cada planta.

Merece la pena subrayar la naturaleza \textbf{interescuela} de la
colaboraci\'on: los desarrolladores pertenec\'ian a la ETSI y el problema
y sus beneficiarios estaban en la ETSIA. Esa distancia, que enriquec\'ia
el proyecto en lo formativo, tuvo tambi\'en un coste: ninguna de las dos
escuelas ten\'ia el proyecto plenamente como propio, y la instalaci\'on
f\'isica --- que deb\'ia montarse en dependencias de la ETSIA con material
adquirido por la ETSIA, pero por personas de la ETSI --- nunca lleg\'o a
ejecutarse.

\section{Prop\'osito de este informe}

El presente documento tiene tres objetivos:

\begin{enumerate}[leftmargin=*]
    \item Reportar el estado actual del proyecto y los logros alcanzados,
    con la honestidad t\'ecnica que la colaboraci\'on institucional exige.

    \item Analizar la sostenibilidad del proyecto en su forma actual,
    identificando las limitaciones estructurales encontradas y aportando
    datos econ\'omicos, organizativos y t\'ecnicos concretos.

    \item Proponer una estrategia de cierre que preserve el valor generado,
    cumpla los objetivos docentes pactados, y mantenga abiertas las
    opciones de continuidad bajo nuevos formatos.
\end{enumerate}

El informe se ha estructurado de modo que sea \'util tanto para una
lectura completa por parte del responsable t\'ecnico como para una
consulta r\'apida por parte de un evaluador acad\'emico: el resumen
ejecutivo y la secci\'on de propuesta de cierre (cap\'itulo~\ref{cap:propuesta})
contienen toda la informaci\'on de decisi\'on necesaria.
